\sectionguard
\section*{Source-Grounded Synthesis Across the Propers}

What follows rests only on the appointed wording, the documented historical
orientation of the texts, and the reception actually checked at its own loci.
It draws no conclusion about why these ten texts stand together, because the
transmission evidence forbids one: the Old Gelasian carries the three orations
as an unassigned Sunday mass, the Hadrianum files them at its own twelfth
Sunday, the Frankish Gelasians at the thirteenth, and the Ottobonianus margins
cue this Sunday's antiphons two masses away. Chants and orations reached each
other by a realignment nobody documented. Everything below is therefore a
statement about the formulary as the 1962 book prints it, never about a
compiler's mind.

\textbf{The formulary's chief coherence is a shared boundary, and it is
verifiable.} Four of its texts stop where a human contribution would begin, and
in three of the four the stopping is a demonstrable editorial act, not the
shape of the source. The Epistle ends mid-verse at \latin{vácua non fuit},
before \latin{sed abundántius illis ómnibus laborávi}. The Communion antiphon
drops the imperative \latin{da ei}, so that the two verbs left standing are
\textit{honour} and \textit{shall be filled}. The Collect's second petition
asks God to add what \latin{orátio} does not presume --- a petition whose
grammatical subject is the prayer, not the person praying. And the Gospel's
protagonist contributes nothing whatever until his tongue has been loosed: he
is brought, taken aside, touched, and opened, and the first thing recorded of
his speech is that it was right. The four boundaries are independently
verified in \texttt{propers/verified.md}; no witness reads them together, and
this synthesis claims only that they are all four there.

\textbf{The reception of five separate elements converges on the
Resurrection, and the convergence is a fact about the commentators rather than
about the Mass.} Hilary, Augustine and Bellarmine read the Introit's
\latin{virtútem et fortitúdinem} of the risen body. Augustine glosses the
Gradual's \latin{reflóruit caro mea} as \latin{resurrexit caro mea}, and
Cassiodorus, Lombard, Aquinas and Bellarmine develop it. Jerome, Augustine,
Cassiodorus, Aquinas, Theodoret and Bellarmine all read the Offertory psalm's
dedication of David's house as the Resurrection. The Epistle is the Church's
oldest resurrection kerygma. Schuster reads the Postcommunion's help to mind
and body as the pledge of the final resurrection. Five elements, six centuries
of witnesses, one direction --- and not one of those witnesses is commenting on
this Sunday. The honest statement is that a reader who arrives with the
tradition in hand will hear a resurrection Mass here, and that the texts
themselves supply that reading only through their reception.

\textbf{The tradition genuinely disagrees at four points, and the formulary
resolves none of them.} At Introit v.~36 the gift of strength is the
resurrection of the body (Hilary, Augustine, Bellarmine), or patience and faith
(Cassiodorus), or providential support (Theodoret). At Gradual v.~7
\latin{reflóruit} is the Resurrection for the whole Latin line and recovered
bodily vigour for Theodoret. At the Alleluia the instruments are allegory for
Augustine and Cassiodorus and instruments for Theodoret and Bellarmine. A
fourth disagreement runs through the Gospel: the deaf-mute is the race in Adam
(ps-Jerome), or any uncatechised soul (Bede), or a particular man with a
demonic affliction (Theophylact), or human nature restored in Christ's own
humanity (the Victor tradition) --- and the sigh is an example, a taking-on of
our cause, a prayer, or a pity, depending on which of them is speaking. These
are not variants of one reading. They yield different sermons, and the guide
prints them as four positions rather than one consensus.

\textbf{Two of the day's texts are the patristic tradition's own language
turned into liturgy, which is a different relation from quotation.} The
Introit sings \latin{unánimes}, which is Augustine's gloss on \latin{unius
modi} and Cassiodorus's lemma; the chant therefore carries into the Mass a word
the Vulgate does not have and the Douay cannot render. The baptismal
\latin{Ephpheta} runs the other way: a word from this Gospel became a rite,
attested by Ambrose at Milan around 390 and still printed in the
\work{Rituale Romanum} of 1872 with the same saliva and the same ears. In both
cases the traffic between exegesis, chant and ritual is documented; in neither
does it license an inference about who put this Mass together.

\textbf{The Epistle's boundary constrains what may be preached from this
Sunday, and the constraint is worth stating positively.} What the Church
actually reads is a creed received and handed on, a list of witnesses, a man
who calls himself untimely born, and the sentence \latin{grátia Dei sum id quod
sum}. What she does not read is the labour and the \latin{mecum}. The lection
is thus, as read, a statement about the origin of a person and a ministry
rather than about the mechanics of cooperation --- and Bede, expounding the
day's Communion antiphon three centuries before the scholastics, says the same
thing about giving: one honours the Lord from one's substance by crediting all
the good one has to grace and not to one's own powers and merits. The Epistle and the
Communion make one claim in two idioms, personal and economic. That is a
source-grounded convergence of two checked expositions, not a claim that either
text was chosen for the other.

\textbf{The orations do the work that the readings leave undone, and they do it
in the plural.} The Collect asks for pardon and for a gift beyond request; the
Secret asks that one offering be at once acceptable to God and support to the
frailty of the offerers; the Postcommunion asks that the sacrament be felt as
help of mind and body, and that, saved in both, we may glory in the fullness of
the remedy. Each is a double request, and in each the second half concerns the
petitioners' weakness. The three prayers travelled together from the Old
Gelasian onward, which is a documented fact about them and the only unity the
formulary can be shown to have been given by anybody. Their shared shape ---
one act of God, two effects, the second of them ours --- is the nearest thing
this Mass has to a thesis, and it is legible without any theory of how the
chants came to join them.
